%% kapa: exam=1 round=23 year=2012 period=1 form=A subject=civil_law subject_ko=민법 seq=21 number=21
%% source: Q-Net 감정평가사 기출문제 (cst00309) artlSeq=5052578 file=raw/1008/5052578/2012년 제23회 감정평가사 제1차시험 기출문제(1교시 A형).pdf
%% conversion: kordoc 4.24.51 -> tools/gen1.py (auto); review=proofread-sonnet; flags=none
\begin{question}{21}
착오에 관한 설명으로 옳지 않은 것은? (다툼이 있으면 판례에 의함)
\begin{choices}
  \item 동기의 착오를 이유로 계약을 취소하기 위해서는 그 동기가 상대방에게 표시되고 의사표시 내용의 중요부분으로 인정될 수 있어야 한다.
  \item 특별한 사정이 없는 한 매매목적물의 시가는 의사표시 내용의 중요부분에 해당하므로 시가에 대한 착오를 이유로 매매계약을 취소할 수 있다.
  \item 착오를 이유로 계약이 취소됨으로써 상대방이 손해를 입었더라도 상대방은 착오자에 대하여 불법행위를 원인으로 그 배상을 청구할 수는 없다.
  \item 착오가 상대방의 기망에 의하여 유발된 경우에는 사기를 이유로 계약을 취소할 수 있다.
  \item 상대방에 의하여 유발된 동기의 착오가 법률행위의 중요부분에 해당하는 경우에는 계약을 취소할 수 있다.
\end{choices}
\end{question}
\answer{2}
